\section{Grafi: Struttura (Lab 20)}

\subsection{Cos'è un Grafo?}
Un Grafo è un insieme di punti (**Nodi** o Vertici) collegati da linee (**Archi** o Edges).
\begin{itemize}
    \item **Orientato:** Le linee sono frecce (A va a B, ma non viceversa).
    \item **Non Orientato:** Le linee sono strade a doppio senso (A è collegato a B).
\end{itemize}

\begin{verbatim}
    (A) -----> (B)
     |          |
     v          v
    (C) <----- (D)
\end{verbatim}

\subsection{Come salvarlo in memoria?}
\begin{concept}{Lista di Adiacenza}
Il modo migliore in JavaScript è usare una \textbf{Mappa} dove:
Chiave = Il Nodo, Valore = Lista dei vicini (Array).
\end{concept}
\begin{jsbox}
class Graph {
    constructor() {
        this.adjList = new Map();
    }

    addNode(nodo) {
        if (!this.adjList.has(nodo)) {
            this.adjList.set(nodo, []); // Creo lista vuota
        }
    }

    addEdge(partenza, arrivo) {
        this.addNode(partenza);
        this.addNode(arrivo);
        
        // Aggiungo il collegamento
        this.adjList.get(partenza).push(arrivo);
        
        // Se fosse NON orientato, dovrei aggiungere anche il viceversa:
        // this.adjList.get(arrivo).push(partenza);
    }
}

let g = new Graph();
g.addEdge("A", "B"); // A -> B
g.addEdge("A", "C"); // A -> C
\end{jsbox}
